% Created 2024-01-08 lun 09:21
% Intended LaTeX compiler: pdflatex
\documentclass[11pt]{article}
\usepackage[utf8]{inputenc}
\usepackage[T1]{fontenc}
\usepackage{graphicx}
\usepackage{longtable}
\usepackage{wrapfig}
\usepackage{rotating}
\usepackage[normalem]{ulem}
\usepackage{amsmath}
\usepackage{amssymb}
\usepackage{capt-of}
\usepackage{hyperref}
\usepackage{minted}

\usepackage{parskip}
\usepackage[utf8]{inputenc} %% For unicode chars
\usepackage{placeins}
\usepackage[margin=2.50cm]{geometry}
\usepackage[T1]{fontenc}
\usepackage{mathpazo}
\linespread{1.05}
\usepackage[scaled]{helvet}
\usepackage{courier}
\hypersetup{colorlinks=true,linkcolor=blue}
\RequirePackage{fancyvrb}
\DefineVerbatimEnvironment{verbatim}{Verbatim}{fontsize=\small,formatcom = {\color[rgb]{0.5,0,0}}}
\usepackage{mdframed}
\BeforeBeginEnvironment{minted}{\begin{mdframed}}
\AfterEndEnvironment{minted}{\end{mdframed}}
\setcounter{secnumdepth}{3}
\date{}
\title{CRUD Car Dealereship OOP}
\hypersetup{
 pdfauthor={},
 pdftitle={CRUD Car Dealereship OOP},
 pdfkeywords={},
 pdfsubject={},
 pdfcreator={Emacs 28.1 (Org mode 9.5.5)}, 
 pdflang={Spanish}}
\begin{document}

\maketitle
\setcounter{tocdepth}{3}
\tableofcontents

\setlength\parindent{10pt}

\section{CRUD Car Dealership in OOP PHP}
\label{sec:org6935730}
Rewriting the Car Dealership PHP CRUD application in an Object-Oriented
Programming (OOP) style involves creating classes to represent entities
like cars, clients, and purchases and encapsulating the CRUD operations
within these classes. Below is a simplified example to demonstrate this
approach.

\subsection{Define the \texttt{Car} Class}
\label{sec:orgaa6fa42}
This class represents a car and includes methods for CRUD operations
related to cars.

\begin{minted}[]{php}
<?=
class Car {
    private $pdo;

    public function __construct($pdo) {
        $this->pdo = $pdo;
    }

    public function create($model, $price, $brand) {
        // Insert a new car into the database
    }

    public function read($id) {
        // Retrieve a car from the database by ID
    }

    public function update($id, $model, $price, $brand) {
        // Update car details in the database
    }

    public function delete($id) {
        // Delete a car from the database
    }
}
\end{minted}

\subsection{Implement CRUD Methods}
\label{sec:org3333f36}
Each method in the \texttt{Car} class should perform the corresponding database
operation. For example, the \texttt{create} method could be implemented as
follows:

\begin{minted}[]{php}
<?=
public function create($model, $price, $brand) {
    $stmt = $this->pdo->prepare("INSERT INTO cars (model, price, brand) VALUES (?, ?, ?)");
    $stmt->execute([$model, $price, $brand]);
}
\end{minted}

Similarly, you would implement the \texttt{read}, \texttt{update}, and \texttt{delete}
methods.

\subsection{Create Other Classes}
\label{sec:org81caa48}
Follow a similar pattern to create classes for other entities like
\texttt{Client} and \texttt{Purchase}. Each class should have its own CRUD methods.

\subsection{Use the Classes in Your Application}
\label{sec:org3b71729}
To use these classes in your application:

\begin{enumerate}
\item \textbf{Establish a Database Connection:} Create a PDO connection to your
database and pass it to the constructor of your entity classes.

\begin{minted}[]{php}
<?=
$pdo = new PDO("mysql:host=localhost;dbname=car_dealership", "username", "password");
\end{minted}

\item \textbf{Perform Operations:} Create instances of your classes and call their
methods to perform operations.

\begin{minted}[]{php}
<?=
 $car = new Car($pdo);
 $car->create("Tesla Model S", 79900, "Tesla");

 $carToUpdate = 1;
 $car->update($carToUpdate, "Tesla Model S Plaid", 129990, "Tesla");

 $carToDelete = 2;
 $car->delete($carToDelete);
\end{minted}
\end{enumerate}

\subsection{Conclusion}
\label{sec:orgf0dda97}
In this OOP approach, each class corresponds to an entity in your
application, and the CRUD operations are methods within these classes.
This structure leads to more organized and maintainable code, especially
for larger applications. Each class is responsible for the operations
related to a single aspect of the application, adhering to the Single
Responsibility Principle, a key tenet of OOP.

\section{Continuation}
\label{sec:orgd1f4fb6}
To complete the Car Dealership CRUD application in an OOP style, we'll
create \texttt{Dealer} and \texttt{Purchase} classes. These classes will handle
operations related to car dealers and purchases, respectively. We'll
also outline the necessary HTML modifications for integrating these
classes into the application.

\subsection{\texttt{Dealer} Class}
\label{sec:orgf537bd9}
This class manages dealer-related operations.

\begin{minted}[]{php}
class Dealer {
    private $pdo;

    public function __construct($pdo) {
        $this->pdo = $pdo;
    }

    public function sellCar($carId, $clientId, $dealerId) {
        // Assuming the method to register a car sale
        // This could include reducing the inventory count for the car
        // and adding an entry to the purchase history

        // Example of adding an entry to a purchase table
        $stmt = $this->pdo->prepare("INSERT INTO purchases (car_id, client_id, dealer_id) VALUES (?, ?, ?)");
        $stmt->execute([$carId, $clientId, $dealerId]);
    }

    // Additional methods for dealer operations can be added here
}
\end{minted}

\subsection{\texttt{Purchase} Class}
\label{sec:org8935ce9}
This class encapsulates operations related to purchases.

\begin{minted}[]{php}
class Purchase {
    private $pdo;

    public function __construct($pdo) {
        $this->pdo = $pdo;
    }

    public function getPurchaseDetails($purchaseId) {
        $stmt = $this->pdo->prepare("SELECT * FROM purchases WHERE id = ?");
        $stmt->execute([$purchaseId]);
        return $stmt->fetch();
    }

    // Additional methods for handling purchase operations
}
\end{minted}

\subsection{HTML Modifications}
\label{sec:orgafbc99f}
To integrate these new classes into your web application, you'll need to
modify the HTML to allow for dealer operations like selling a car.
Here's an example of how you might structure a simple form for a dealer
to record a car sale:

\subsubsection{HTML Form for Selling a Car (\texttt{sell\_car.html})}
\label{sec:orgc1a2d11}
\begin{minted}[]{html}
<!DOCTYPE html>
<html>
<head>
    <title>Sell a Car</title>
</head>
<body>
    <h1>Sell a Car</h1>
    <form action="sell_car.php" method="post">
        Car ID: <input type="number" name="car_id"><br>
        Client ID: <input type="number" name="client_id"><br>
        Dealer ID: <input type="number" name="dealer_id"><br>
        <input type="submit" value="Sell Car">
    </form>
</body>
</html>
\end{minted}

\subsubsection{PHP Script for Handling the Form (\texttt{sell\_car.php})}
\label{sec:org5de82da}
\begin{minted}[]{php}
<?php
require_once 'db_connect.php'; // Database connection
require_once 'Dealer.php';

$pdo = new PDO("mysql:host=localhost;dbname=car_dealership", "username", "password");

if ($_SERVER["REQUEST_METHOD"] == "POST") {
    $dealer = new Dealer($pdo);
    $dealer->sellCar($_POST['car_id'], $_POST['client_id'], $_POST['dealer_id']);

    echo "Car sold successfully.";
}
?>
\end{minted}

\subsection{Conclusion}
\label{sec:orge6bff31}
In the OOP version of your Car Dealership application, \texttt{Dealer} and
\texttt{Purchase} classes encapsulate all the logic related to their respective
domains. This makes your code more organized, modular, and easier to
maintain.

You'll need to adjust your HTML forms to work with these classes. The
example above shows how you can create a form for a dealer to sell a
car, which is then processed by a PHP script using the \texttt{Dealer} class.

Remember, these are basic examples to demonstrate the concept. In a
real-world application, you'd likely have more comprehensive error
checking, input validation, and possibly more complex business logic
incorporated into your class methods.
\end{document}
